\chapter{Conclusiones}

Este capítulo presenta las conclusiones del trabajo de grado organizadas según los cuatro objetivos específicos planteados en la Sección~2.3. Cada apartado sintetiza las actividades ejecutadas y las cifras experimentales verificadas en los capítulos previos (Capítulos~4 a~7). Adicionalmente, se incluye una conclusión metodológica sobre el efecto de la unidad de partición en la evaluación de modelos dermatoscópicos.

\begin{itemize}
  \item \textbf{Gestión, preprocesamiento y aumento de la base de datos (Objetivo específico 1):} se procesaron las 10.015 imágenes dermatoscópicas correspondientes a 7.470 lesiones de las siete clases del repositorio HAM10000~\cite{tschandl2018ham10000} (Tabla~\ref{tab:distribucion}). Mediante \texttt{StratifiedGroupKFold}~\cite{pedregosa2011sklearn}, los datos se dividieron por \texttt{lesion\_id} en 7.153 imágenes de entrenamiento, 1.431 de validación y 1.431 de prueba sin lesiones compartidas entre subconjuntos (Tabla~\ref{tab:particion}). A todas las imágenes se les aplicó redimensionado con relleno negro, filtro de mediana $3 \times 3$, remoción de vello~\cite{lee1997dullrazor,telea2004inpainting} y ecualización CLAHE con límite 2,0~\cite{zuiderveld1994clahe}. El desbalance se trató en el entrenamiento mediante aumento geométrico y fotométrico en línea junto con pesos de clase entre 0,46 y 3,53 (Tabla~\ref{tab:pesos}).

  \item \textbf{Implementación de arquitecturas basadas en redes neuronales convolucionales (Objetivo específico 2):} se implementaron siete arquitecturas de clasificación de imágenes dermatoscópicas evaluadas en dos resoluciones de entrada de 224\,px y 384\,px (Tabla~\ref{tab:arquitecturas}). El conjunto de modelos incluyó cinco redes preentrenadas en ImageNet~\cite{deng2009imagenet} (DenseNet121~\cite{huang2017densenet}, ResNet50~\cite{he2016resnet}, MobileNetV3Large~\cite{howard2019mobilenetv3}, EfficientNetV2-S~\cite{tan2021efficientnetv2} y ConvNeXt-Tiny~\cite{liu2022convnext}), una red híbrida DenseNet121-BiLSTM y una red CNN-BiLSTM entrenada desde cero. Las redes con transferencia de aprendizaje se ajustaron con una receta de dos fases (hasta 15 épocas con base congelada y hasta 40 épocas de ajuste fino con decaimiento coseno). Ocho configuraciones se entrenaron con tres semillas aleatorias (42, 43 y 44) y se compararon en el conjunto de prueba de 1.431 imágenes (Tabla~\ref{tab:resultados-modelos}).

  \item \textbf{Evaluación y selección del modelo para identificación de melanoma (Objetivo específico 3):} mediante la regla de selección preregistrada en validación ($\delta = 0{,}0112$ y desempate por AUC de melanoma), se eligió MobileNetV3Large a 384\,px con semilla 44~\cite{howard2019mobilenetv3}. En la prueba de HAM10000 con la clase más probable, este modelo obtuvo 81,8\,\% de accuracy, 0,647 de Macro-F1 y 0,903 de AUC de melanoma (Tabla~\ref{tab:modos-decision}). En la validación sobre las 1.512 imágenes de prueba de ISIC 2018~\cite{codella2019isic2018}, el modelo alcanzó un AUC de melanoma de 0,908 y un Macro-F1 de 67,9\,\% sin modificar sus parámetros (Tabla~\ref{tab:validacion-isic}). Con la calibración por temperatura ($T = 3{,}2933$)~\cite{guo2017calibration}, la sensibilidad para melanoma en HAM10000 alcanzó 72,3\,\% en el modo equilibrado ($t = 0{,}20$) y 80,5\,\% en alta sensibilidad ($t = 0{,}11$).

  \item \textbf{Integración del modelo seleccionado en el aplicativo web (Objetivo específico 4):} el modelo seleccionado se exportó al formato ONNX (12,9\,MB) y se integró en un aplicativo web estático desarrollado con React, TypeScript y Vite (Sección~6.1). El aplicativo ejecuta en un Web Worker el preprocesamiento con OpenCV.js 5.0.0 y la inferencia con ONNX Runtime Web 1.31.1~\cite{ning2024onnxweb} sin transmitir imágenes fuera del navegador. En las 120 imágenes de prueba de paridad, el aplicativo coincidió en el 100\,\% de las decisiones con la cadena de referencia en Python con \texttt{onnxruntime} (Sección~6.6). El procesamiento tomó en promedio 59\,ms por imagen en Chromium 158, se verificó con cero violaciones automáticas en \texttt{axe-core} y se desplegó como sitio web estático.

  \item \textbf{Hallazgo metodológico sobre la unidad de partición:} la comparación entre esquemas de división en HAM10000 mostró que partir por imagen deja el 36,1\,\% de la prueba con otra fotografía de la misma lesión en entrenamiento (Sección~5.5). En la clase \texttt{mel}, esta superposición entre entrenamiento y prueba afectó al 63,5\,\% de los melanomas evaluados (106 de 167 imágenes). Como resultado, la partición por imagen infló el Macro-F1 de prueba cerca de 9 puntos porcentuales ($+8{,}7$ puntos en DenseNet121 y $+9{,}2$ puntos en CNN-BiLSTM, Tabla~\ref{tab:comparacion-particiones})~\cite{kapoor2023leakage,abhishek2025quality}.
\end{itemize}
